% ============================================================
% macros/nonvf.tex -- chapter-local macros for ch:nonvf.
% Merged from nonvf-paper/preamble.sty (long version, lines
% 104-191) and the inline definitions in nonvf-paper-succint/
% main.tex. Loaded inside a group, so the overrides below vanish
% after the chapter.
%
% Macro decisions:
%   * \bools is overridden to {0,1} (succinct paper's convention;
%     the thesis-global is {\top,\bot}).
%   * \Prob: the succinct paper defines \Prob[opt]{arg} =
%     \mathbf{Pr}^{opt}[arg]; the long paper's \Prob is a bare
%     \mathbf{Pr}. The thesis-wide symbol is \mathbf{P}, so the
%     chapter uses ONE consistent rendering: \Prob[opt]{arg} =
%     \mathbf{P}^{opt}[arg] (succinct calling convention, thesis
%     symbol). The long version's bare-\Prob usages lifted into
%     the chapter have been normalized to \Prob{...}; the bare
%     symbol is available as \Problem (as in the succinct paper).
%   * \toc: the long paper's \toc = \to_{\mathsf{c}} (partial
%     computable function) is kept, since it is used throughout
%     the program-market construction. The succinct paper's
%     conflicting inline \toc = \to_{\Compcl} was used once, in
%     the Discussion; that occurrence is written literally as
%     \to_{\Compcl} in the chapter.
%   * \Nats \Rats \Reals \Bools \probs \Props \strings \zero
%     \finset \supp \Expect and the theorem environments
%     (incl. lefinition, hproof) are already global -- not
%     redefined here.
% ============================================================

% --- overrides of thesis-global macros ---
\DeclareDocumentCommand{\bools}{}{\{0,1\}}
\DeclareDocumentCommand{\Problem}{}{\mathbf{P}}
\DeclareDocumentCommand{\Prob}{O{}m}{%
  \mathbf{P}%
  \IfValueT{#1}{^{#1}}%
  \left[#2\right]%
}

% --- basic math (long version, preamble.sty) ---
\NewDocumentCommand{\code}{m}{\left\langle#1\right\rangle}
\NewDocumentCommand{\concat}{m m}{#1:#2}
\NewDocumentCommand{\compl}{}{\,\raisebox{0.25ex}{\rotatebox[origin=c]{180}{\ensuremath{\DOTSB\lhook\joinrel\relbar}}}}
\NewDocumentCommand{\proj}{O{}}{\operatorname{pr}_{#1}}
\NewDocumentCommand{\too}{}{\to_\circ}
\NewDocumentCommand{\topt}{}{\to_{\mathsf{p}}}
\NewDocumentCommand{\toc}{}{\to_{\mathsf{c}}}
\NewDocumentCommand{\mapstopt}{}{\mapsto_{\mathsf{p}}}
\NewDocumentCommand{\mapstoc}{}{\mapsto_{\mathsf{c}}}
\NewDocumentCommand{\toe}{}{\to_{-}}
\NewDocumentCommand{\mapstoe}{}{\mapsto_{-}}
\NewDocumentCommand{\tod}{}{\to_{\downarrow}}
\NewDocumentCommand{\mapstod}{}{\mapsto_{\downarrow}}
\NewDocumentCommand{\tode}{}{\to_{-\downarrow}}
\NewDocumentCommand{\mapstode}{}{\mapsto_{-\downarrow}}
\NewDocumentCommand{\intlen}{m}{\left|#1\right|}
\NewDocumentCommand{\maybe}{m}{\overline{#1}}
\NewDocumentCommand{\preimg}{m m}{{#1}^{-1}\left(#2\right)}

% --- logic (long version) ---
\NewDocumentCommand{\play}{m m}{#1[#2]}
\NewDocumentCommand{\moves}{}{\maybe{\finset{\Nats}}}
\NewDocumentCommand{\playfull}{}{\operatorname{Game}}
\NewDocumentCommand{\halts}{}{\operatorname{Halts}}

% --- agent-notation system (long version, LaTeX3) ---
\ExplSyntaxOn
% Guarded so this file can be loaded more than once (it is pre-loaded
% globally for the ToC/LoF/LoT pass and again inside the chapter group).
\prop_if_exist:NF \g_based_agents_prop { \prop_new:N \g_based_agents_prop }
\prop_gset_from_keyval:Nn \g_based_agents_prop
  {
    a = \alpha, b = \beta, c = \gamma, m = \mu, p = \pi, pp = \varpi, s = \mathcal{A}
  }

\NewDocumentCommand{\agent}{O{a}}{
  \prop_get:NnNTF \g_based_agents_prop {#1} \l_tmpa_tl {
    \tl_use:N \l_tmpa_tl
  }{
    #1
  }
}
\ExplSyntaxOff

\NewDocumentCommand{\trader}{O{a}}{\hat{\agent[#1]}}
\NewDocumentCommand{\player}{O{a}}{\tilde{\agent[#1]}}
\NewDocumentCommand{\instaplayer}{O{a}}{\breve{\agent[#1]}}
\NewDocumentCommand{\labeler}{O{a}}{{\agent[#1]}^{\#}}
\NewDocumentCommand{\endowment}{O{a}}{{\agent[#1]}^{\operatorname{b}}}
\NewDocumentCommand{\birthday}{O{a}}{{\agent[#1]}^{\operatorname{t}}}
\NewDocumentCommand{\inventory}{O{a}}{{\agent[#1]}^{\$}}
\NewDocumentCommand{\change}{O{} O{a}}{{\delta\agent[#2]}^{\$}_{#1}}
\NewDocumentCommand{\father}{O{m}}{{\agent[#1]}^{\operatorname{bh}}}
\NewDocumentCommand{\price}{O{p}}{{\agent[#1]}^{*}}
\NewDocumentCommand{\longrun}{O{pp}}{{\agent[#1]}^{\infty}}
\NewDocumentCommand{\tpler}{O{s}}{{\agent[#1]}^{\circ}}
\NewDocumentCommand{\foo}{}{\texttt{``foo''}}

% special stupid thing because writing \player[s] \instaplayer[s] in theorem name confuses LaTex
\NewDocumentCommand{\instaplayers}{}{\instaplayer[s]}
\NewDocumentCommand{\players}{}{\player[s]}

% --- prediction markets (long version) ---
\NewDocumentCommand{\finprocs}{}{\operatorname{FinProc}}
\NewDocumentCommand{\portfolios}{}{\operatorname{PF}}
\NewDocumentCommand{\pricee}{}{\pi}
\NewDocumentCommand{\inventorye}{}{{\boldsymbol{\alpha}}}

% --- Garrabrant induction appendix (long version) ---
\NewDocumentCommand{\thms}{}{\Theta}
\NewDocumentCommand{\world}{}{w}
\NewDocumentCommand{\PC}{}{\mathsf{PC}}
\NewDocumentCommand{\indicate}{m}{\mathbf{I}\left[#1\right]}

% --- inline definitions from the succinct version's main.tex ---
\NewDocumentCommand{\Ack}{}{\mathcal{A}}
\NewDocumentCommand{\Compcl}{}{\mathcal{C}}
\newcommand{\noop}[1]{}
